\section{Модель системи}

\subsection{Діаграма прецедентів}

Систему використовують чотири категорії акторів (рис.~\ref{fig:use_case.png}): користувач
веб-інтерфейсу, користувач CLI-клієнта, дослідник, який запускає навантажувальні тести, і
оркестратор Docker, що перевіряє стан контейнера. Зовнішня система Ollama бере участь у
розширенні запиту та в перевірці стану. Прецеденти відповідають ендпоінтам фіксованого
контракту: побудова індексу (UC1), пошук через REST (UC2), пошук зі стрімінгом етапів через
WebSocket (UC3), перегляд статистики (UC4) і перевірка стану (UC5). Розширення запиту
включене в UC3 завжди (його можна вимкнути полем \code{expand:false}), а в UC2 розширює
його лише за явного прапорця \code{expand:true} -- за замовчуванням REST-пошук працює у
чистому режимі \code{keyword}.

\fig[0.98\textwidth]{use_case.png}{Діаграма прецедентів}

\subsection{Діаграма класів}

Код поділено на кореневий пакет і вісім підпакетів (рис.~\ref{fig:class_overview.png}). Пакет \code{core} містить
обчислювальне ядро: власний пул потоків \code{ThreadPool}, оркестратор побудови
\code{Indexer}, інтерфейс \code{IndexBuilder} з двома реалізаціями спільної структури
(\code{PartitionedIndexBuilder} -- неблокуюча, \code{StripedIndexBuilder} -- блокувальна),
незмінний індекс \code{InvertedIndex} з таблицею документів \code{DocumentTable},
неблокуючий стек \code{TreiberStack}, список входжень \code{PostingList}, токенізатор і
байтову таблицю термів задачі \code{BatchTermTable}. Пакет \code{net} реалізує мережевий
рівень (\code{HttpServer}, \code{HttpRequest}, \code{HttpResponse}, \code{WebSocket}), пакет
\code{api} -- маршрутизацію і спільний для REST і WebSocket конвеєр пошуку
(\code{Router}, \code{SearchService}, \code{WsSearchHandler}), пакет \code{search} --
ранжування BM25 і сніпети, пакет \code{llm} -- клієнт Ollama і розширювач запитів. Класи
\code{SearchCli} і \code{Bench} є окремими точками входу того самого jar-файла. Детальні
діаграми класів із полями синхронізації наведено в додатку~А (рис.~А.1--А.2).

\fig[0.98\textwidth]{class_overview.png}{Діаграма класів (огляд за пакетами)}

\subsection{Діаграми послідовності}

\textit{Побудова індексу} (рис.~\ref{fig:seq_index.png}). \code{Router} перевіряє параметри
\code{threads} (1--256), \code{partitions} (1--4096), \code{batch\_size} і назву реалізації;
ціле число приймається також у вигляді \code{4.0} або \code{"4"}, інші значення дають 400.
\code{Indexer} захоплює \code{buildLock} методом \code{tryLock}: паралельний запит на
перебудову не стає в чергу, а одразу отримує 409. Далі на викликаючому потоці формується
відсортований список файлів -- позиція файла в ньому є ідентифікатором документа, тому
результат не залежить від кількості потоків. Список ріжеться на батчі по 32 документи, і
кожен батч стає окремою задачею пулу. Задача читає й токенізує свої документи у приватну
таблицю без синхронізації і передає її будівнику; будівник розбиває таблицю на фрагменти за
партиціями і зливає їх неблокуючим протоколом (підрозділ~2.5). Після завершення всіх задач
по одній задачі на партицію сортує і «заморожує» списки входжень, а новий знімок індексу
атомарно підміняє попередній.

\fig{seq_index.png}{Діаграма послідовності: побудова індексу}

\textit{Пошук через REST} (рис.~\ref{fig:seq_search_rest.png}). Некоректний JSON, порожній
запит або непідтримуваний режим дають 400. За прапорця \code{expand} запит спершу
розширюється: результат береться з LRU-кешу, або з уже виконуваного однакового запиту, або
з нового виклику моделі в межах семафора і дедлайну; помилка чи тайм-аут моделі дають
порожній список термів, і пошук продовжується. Знімок індексу береться вже \textit{після}
розширення, тому запит, що чекає на модель, не утримує в пам'яті індекс, який паралельна
перебудова вже замінила. Ранжування BM25 і формування сніпетів виконуються без блокувань.

\fig{seq_search_rest.png}{Діаграма послідовності: пошук через REST API}

\textit{Пошук через WebSocket} (рис.~\ref{fig:seq_search_ws.png}). Після рукостискання сокет
передається з пулу HTTP-потоків у пул WebSocket-сеансів. В одному сеансі можна надіслати
кілька запитів поспіль; на кожен сервер відповідає трьома повідомленнями-етапами, а для
помилок окремого запиту (некоректний JSON, індекс ще не побудовано) надсилає повідомлення з
етапом \code{error}, не закриваючи з'єднання. Порушення протоколу закривають сеанс із кодом
RFC~6455.

\fig{seq_search_ws.png}{Діаграма послідовності: пошук через WebSocket}

\textit{Статистика та перевірка стану} (рис.~\ref{fig:seq_stats.png},
\ref{fig:seq_health.png}). Обидва запити лише читають знімок індексу через volatile-посилання.
Перевірка стану зондує Ollama запитом \code{GET /api/tags} з тайм-аутом 1~с і кешує
результат на 2~с, щоб часті звернення healthcheck не створювали навантаження на модель.

\fig[0.85\textwidth]{seq_stats.png}{Діаграма послідовності: статистика}
\fig[0.85\textwidth]{seq_health.png}{Діаграма послідовності: перевірка стану}

\subsection{Діаграма розгортання}

Система розгортається на одній машині (рис.~\ref{fig:deployment.png}) у проєкті Docker
Compose \code{pocw-search}. Сервіс \code{server} (образ на базі \code{eclipse-temurin:21-jre},
порт 8080) містить jar-файл, статичний Web UI та обгортки \code{search} і \code{bench}.
Сервіс \code{ollama} (порт 11434) зберігає ваги моделі в іменованому томі \code{ollama\_data},
тому повторні запуски не завантажують модель знову. Два одноразові сервіси виконуються під
час першого запуску: \code{ollama-init} завантажує модель, \code{dataset-init} розрізає архів
корпусу на документи в іменованому томі \code{corpus}, який сервер монтує тільки для
читання. Сервер стартує після завершення \code{dataset-init}, але не чекає на модель: до її
завантаження пошук працює без розширення.

\fig[0.82\textwidth]{deployment.png}{Діаграма розгортання (топологія Docker Compose)}
